\documentclass[10pt,a4paper]{amsart}
\usepackage[margin=19mm]{geometry}
\usepackage{amsmath,amssymb,mathtools}
\usepackage[scr=rsfs,cal=euler]{mathalfa}
\usepackage{xfrac}
\usepackage[unicode,hidelinks,bookmarks=false]{hyperref}
\newcommand{\ie}{i.e.\ }
\newcommand{\cf}{cf.\ }
\newcommand{\resp}{respectively\ }
\newcommand{\Subsection}{\subsection}
\providecommand{\nobreakdash}{\nobreak\hbox{-}}
\setlength{\emergencystretch}{2em}
\multlinegap0pt
\usepackage{bm}
\usepackage{bbm}
\usepackage{dsfont}
\usepackage{xspace}
\usepackage{cancel}
\usepackage{mathtools}
\usepackage{amsthm}
\makeatletter
\@ifundefined{theorem}{\newtheorem{theorem}{Theorem}}{}
\makeatother
\newcommand{\Hy}{\mathcal{H}}
\newcommand{\I}{\mathcal{I}}
\newcommand{\J}{\mathcal{J}}
\newcommand{\subs}{\mathcal{C}}
\title[Multiple testing]{Multiple testing}
\author{Jesse Hemerik}
\date{Source: arXiv:2606.26781v1 (CC BY 4.0)}
\begin{document}
\maketitle

\noindent\emph{Source and licence.} Multiple testing. Version arXiv:2606.26781v1 (CC BY 4.0), distributed under the Creative Commons Attribution 4.0 International licence, as recorded in the arXiv licence metadata for this exact version and shown on the abstract page https://arxiv.org/abs/2606.26781v1. Full rights notice: licenses/arxiv-2606.26781-CC-BY-4.0.txt.\par\smallskip

\noindent\textit{Editorial scope.} Counted unit: the Theorem whose source label is \texttt{thmHolmisCTP} in the article (source0.tex lines 2059--2061, source label \texttt{thmHolmisCTP}), delivered together with the article's own complete original proof of it (source0.tex lines 2062--2099, one \texttt{proof} environment of 38 lines). No mathematics is written, repaired or re-derived: statement and proof are the article's own text line for line. Required context retained uncounted: localtestHolm (lines 2052--2054). Every definition, display and label of those lines is retained unchanged. Omitted from this package and not redistributed: 1--2051, 2055--2058, 2100--3863. Nothing inside a retained excerpt is omitted or altered: every retained body is delivered exactly as the article's lines are written, so no omission receipt of any kind accompanies this package. Citations: the retained excerpts carry no citation command at all, so no bibliography entry of the article is transcribed and no reference is redistributed. Reference adaptation: none was needed and none was made; every reference inside the delivered unit resolves inside it, so no unresolved reference or citation is produced. Printed slips kept literal: no printed word, symbol, hypothesis or number of the counted statement or of its proof was altered. Layout and class: the article's own class is \texttt{article} with packages including amsmath, amsfonts, amssymb, mathrsfs, bm, bbm, amsthm, comment, tocbibind, graphicx, subcaption, changepage, natbib, abstract; the wrapper is the lane's pinned \texttt{amsart} editorial wrapper at 10pt on A4, which restates the theorem-like environments the retained text needs and adds the article's own macro definitions that the retained text needs (\texttt{\textbackslash Hy}, \texttt{\textbackslash I}, \texttt{\textbackslash J}, \texttt{\textbackslash subs}), with extra wrapper packages (bm, bbm, dsfont, xspace, cancel, mathtools). The delivered numbering is the unit's own. Assets: no retained span contains a figure environment, a graphics-inclusion command, a PDF-page inclusion, a table or a TikZ picture; no image file is copied into this package, the package declares no asset dependency and no image file is redistributed with it. Licence: version arXiv:2606.26781v1 of this article is distributed under the Creative Commons Attribution 4.0 International licence, as recorded in the arXiv licence metadata for this exact version and shown on the abstract page https://arxiv.org/abs/2606.26781v1. A full rights notice is delivered at licenses/arxiv-2606.26781-CC-BY-4.0.txt. Characters: every retained excerpt is pure ASCII, so the delivered engine, XeLaTeX with its default Latin Modern text font, reports no missing character.
\begin{equation} \label{localtestHolm}
\phi_{\I} = \mathbbm{1}\big(\min\{P_i: i\in \I\}\leq \alpha /|\I|\big).
\end{equation}   

\begin{theorem} \label{thmHolmisCTP}
The  CTP with local tests  \eqref{localtestHolm} rejects exactly the same elementary hypotheses as Holm's method.
\end{theorem}

\begin{proof}
Without loss of generality, we can rename the hypotheses such that the  p-values are ordered: $p_{1}\leq ... \leq p_{m}$. (Thus, $\Hy_1$ is the hypothesis with the smallest p-value, etc.) Indeed, we can do this since clearly,  both methods simply reject a set of hypotheses corresponding to the smallest so-many p-values, and how many hypotheses they reject does not depend on the original order.

Note that  the methods reject at least one hypothesis if and only if Bonferroni's global test rejects $\Hy_{\{1,...,m\}}$.
Consequently,  the  CTP rejects at least $\Hy_{1}$ if and only if   Holm rejects at least $\Hy_{1}$.
\\
\\
\emph{Part 1: Showing that  Holm rejects at least as many as the CTP.}
Suppose that the CTP rejects $2\leq k \leq m$ hypotheses. We show that Holm does as well.
%Recall our notation: $\psi_{\{k\}}=1$ indicates whether if  $\Hy_{k}$ is rejected by the CTP.
We give a proof by induction.
Let $1\leq j< k$ and suppose we have shown that Holm rejects $\Hy_1,...,\Hy_j$. We show that Holm also rejects $\Hy_{j+1}$.
Since Holm rejects  $\Hy_1,...,\Hy_j$, we know that $\forall 1\leq i\leq j:  p_i\leq \alpha/(m+1-i)$. We must show that also 
$p_{j+1}\leq \alpha/(m+1-(j+1))$. We know that the CTP rejects $\Hy_{j+1}$, i.e., 
$$ \forall \{{j+1}\} \subseteq \J \in \subs: \phi_{\J}=1  ,$$
i.e., 
$$ \forall\{{j+1}\} \subseteq \J \in \subs: \min\{p_i: i\in \J\}\leq \alpha/ |\J|.$$
Hence, in particular, for $\J=\{j+1,...,m\}$,
$$\min\{p_i: i\in \J\}\leq \alpha/ |\J|,$$
which means that
$$p_{j+1}\leq \alpha/ (m-j),$$ which means that indeed Holm rejects $\Hy_{j+1}$ as well.
\\
\\
\emph{Part 2: Showing that  the CTP rejects at least as many as Holm.}
Suppose that Holm rejects $2\leq k \leq m$ hypotheses. We show that the CTP does as well. We again use induction.
%To show that the CTP  also rejects at least $k$ hypotheses, we proceed by induction.
Let $1\leq j< k$ and suppose we have shown that the CTP  rejects $\Hy_1,...,\Hy_j$. We show that the CTP also rejects $\Hy_{j+1}$.
We know that 
$\forall 1\leq i\leq j+1:  p_i\leq \alpha/(m+1-i)$.
Further, we know that the CTP rejects $\Hy_j$, i.e.,
%\begin{equation} \label{holmproof1}
 $$\forall\{{j}\} \subseteq \J \in \subs: \min\{p_i: i\in \J\}\leq \alpha/ |\J|,$$
% \end{equation}
and we must show that 
$$ \forall\{{j+1}\} \subseteq \J \in \subs: \min\{p_i: i\in \J\}\leq \alpha/ |\J|.$$
Since the CTP rejects $\Hy_1,...,\Hy_j$, we know that once we include an index $l<j+1$ in $\J$, then $\psi_{\{l\}}=1$ and hence  $\phi_{\J}=1$ and hence  $\min\{p_i: i\in \J\}\leq \alpha/ |\J|$.
Hence, we only need to check that for all $\J\subseteq\{j+1,...,m\}$ that contain $j+1$, we have $\min\{p_i: i\in \J\}\leq \alpha/ |\J|$. But this clearly holds since for such $\J$ we have  $\min\{p_i: i\in \J\}=p_{j+1}\leq \alpha/(m+1-(j+1))\leq \alpha/|\J|$, since $\J\leq m-j$ for such $\J$.
\end{proof}

\end{document}
